\documentclass{article}

\usepackage{fancyhdr}
\usepackage{extramarks}

%amsmath
\usepackage{amsmath}
\usepackage{amsthm}
\usepackage{amsfonts}
\usepackage{mathtools}
\usepackage{amssymb}

%Enumerating items and adding columns
\usepackage{enumitem}
\usepackage{multicol}

%TikZ picture
\usepackage{tikz}
\usepackage{scalerel}
\usepackage{pict2e}
\usepackage{tkz-euclide}
\usetikzlibrary{calc}
\usetikzlibrary{patterns,arrows.meta}
\usetikzlibrary{shadows}
\usetikzlibrary{external}
\usetikzlibrary{automata,positioning}

%pgfplots
\usepackage{pgfplots}
\pgfplotsset{compat=newest}
\usepgfplotslibrary{statistics}
\usepgfplotslibrary{fillbetween}

%colours
\usepackage{xcolor}

%
% Basic Document Settings
%

\topmargin=-0.45in
\evensidemargin=0in
\oddsidemargin=0in
\textwidth=6.5in
\textheight=9.0in
\headsep=0.25in

\linespread{1.1}

\pagestyle{fancy}
\lhead{\hmwkAuthorName}
\chead{\hmwkChapter . \hmwkChapterTitle \hmwkClassTime: \hmwkTitle}
\rhead{\firstxmark}
\lfoot{\lastxmark}
\cfoot{\thepage}

\renewcommand\headrulewidth{0.4pt}
\renewcommand\footrulewidth{0.4pt}

\setlength\parindent{0pt}

%
% Create Problem Sections
%

\newcommand{\enterProblemHeader}[1]{
    \nobreak\extramarks{}{Problem \arabic{#1} continued on next page\ldots}\nobreak{}
    \nobreak\extramarks{Problem \arabic{#1} (continued)}{Problem \arabic{#1} continued on next page\ldots}\nobreak{}
}

\newcommand{\exitProblemHeader}[1]{
    \nobreak\extramarks{Problem \arabic{#1} (continued)}{Problem \arabic{#1} continued on next page\ldots}\nobreak{}
    \stepcounter{#1}
    \nobreak\extramarks{Problem \arabic{#1}}{}\nobreak{}
}

\setcounter{secnumdepth}{0}
\newcounter{partCounter}
\newcounter{homeworkProblemCounter}
\setcounter{homeworkProblemCounter}{1}
\nobreak\extramarks{Problem \arabic{homeworkProblemCounter}}{}\nobreak{}

%
% Homework Problem Environment
%
% This environment takes an optional argument. When given, it will adjust the
% problem counter. This is useful for when the problems given for your
% assignment aren't sequential. See the last 3 problems of this template for an
% example.
%
\newenvironment{homeworkProblem}[1][-1]{
    \ifnum#1>0
        \setcounter{homeworkProblemCounter}{#1}
    \fi
    \section{Problem \arabic{homeworkProblemCounter}}
    \setcounter{partCounter}{1}
    \enterProblemHeader{homeworkProblemCounter}
}{
    \exitProblemHeader{homeworkProblemCounter}
}

%
% Homework Details
%   - Title
%   - Due date
%   - Class
%   - Section/Time
%   - Instructor
%   - Author
%

\newcommand{\hmwkTitle}{Exercises}
\newcommand{\hmwkChapter}{CHAPTER 2}
\newcommand{\hmwkDate}{April 3, 2026}
\newcommand{\hmwkClassTime}{}
\newcommand{\hmwkChapterTitle}{Direct Proofs}
\newcommand{\hmwkAuthorName}{\textbf{Christian Jelo R. Artoza}}

%
% Title Page
%

\title{
    \vspace{2in}
    \textmd{\textbf{\hmwkChapter:\ \hmwkTitle}}\\
    \normalsize\vspace{0.1in}\small{\hmwkDate}\\
    \vspace{0.1in}\large{\textit{\hmwkChapterTitle\ \hmwkClassTime}}
    \vspace{3in}
}

\author{\hmwkAuthorName}
\date{}

\renewcommand{\part}[1]{\textbf{\large Part \Alph{partCounter}}\stepcounter{partCounter}\\}

%
% Various Helper Commands
%

% qedsymbol that is always flushed to the right
\newcommand{\qedright}{
    \begin{flushright}
        \qedsymbol
    \end{flushright}
}

% For propositions
\newtheorem{prop}{Proposition}

% Useful for algorithms
\newcommand{\alg}[1]{\textsc{\bfseries \footnotesize #1}}

% For derivatives
\newcommand{\deriv}[1]{\frac{\mathrm{d}}{\mathrm{d}x} (#1)}

% For partial derivatives
\newcommand{\pderiv}[2]{\frac{\partial}{\partial #1} (#2)}

% Integral dx
\newcommand{\dx}{\mathrm{d}x}

% Alias for the Solution section header
\newcommand{\solution}{\textbf{\large Solution}}

% Probability commands: Expectation, Variance, Covariance, Bias
\newcommand{\E}{\mathrm{E}}
\newcommand{\Var}{\mathrm{Var}}
\newcommand{\Cov}{\mathrm{Cov}}
\newcommand{\Bias}{\mathrm{Bias}}

\begin{document}

\maketitle

\pagebreak

%Plot Settings
\pgfplotsset{
    standard/.style={
    axis line style = thick,
    axis equal,
    trig format=rad,
    enlargelimits,
    axis x line=middle,
    axis y line=middle,
    axis line style={latex-latex},
    enlarge x limits=0.15,
    enlarge y limits=0.15,
    every axis x label/.style={at={(current axis.right of origin)},anchor=north west},
    every axis y label/.style={at={(current axis.above origin)},anchor=south east}
    }
}

\begin{homeworkProblem}
    \textit{(Exercise 2.6.)} Determine conditions on integers $m$ and $n$ for which $mn$ is
    even. Write down your conditions as a conjecture, and then prove that your conjecture is
    correct.
    \vspace{1.5mm}

    \textbf{Solution.}

    \begin{prop}
        Let $m$ and $n$ be integers. If either $m$ or $n$ is even, then $mn$ is even.
    \end{prop}

    \begin{proof}
        Without loss of generality, suppose $m$ is even. Then, $m=2k$ for some integer $k$.
        Hence,
        \begin{align*}
            mn=(2k)n=2kn
        \end{align*}
        Since $kn$ is an integer, it follows that $mn$ is even. This completes the proof.
    \end{proof}
\end{homeworkProblem}

\begin{homeworkProblem}
    \textit{(Exercise 2.7.)} Prove the following. For each, $m,n,$ and $t$ are integers.
    \begin{multicols}{2}
        \begin{enumerate}[label=(\alph*)]
            \item If $m \mid n$, then $m^2 \mid n^2$.
            \item If $m \mid n$, then $m \mid (7n^3+12n^2-n)$.
            \item If $m \mid n$ and $m \mid t$, then $m \mid (n+t)$.
            \columnbreak
            \item If $3 \mid 2n$, then $3\mid n$.
            \item If $m^3 \mid n$ and $n^4 \mid t$, then $m^12 \mid t$.
        \end{enumerate}
    \end{multicols}
    \textbf{Solution.}
    \begin{enumerate}[label=(\alph*)]
        \item \begin{proof}
            Assume $m \mid n$. Then, $n = k\cdot m$ for some integer $k$. Multiplying both sides
            of this equation by $m$ we get
            \begin{align*}
                mn = k\cdot m^2
            \end{align*}
            This implies that $m^2 \mid mn$. Now, going back to the equation $n = k \cdot m$, we can
            multiply both sides by $n$ and get
            \begin{align*}
                n^2 = k \cdot mn
            \end{align*}
            Hence, $mn \mid n^2$. So, since $m^2 \mid mn$, it follows by transitivity that $m^2 \mid n^2$,
            as desired.
            \end{proof}
        \item \begin{proof}
            Assume $m \mid n$. Then, $n = k \cdot m$ for some $k$, and hence
            \begin{align*}
                7n^3+12n^2-n
                &= 7(k\cdot m)^3 + 12(k\cdot m)^2 - k\cdot m \\
                &= 7k^3 m^3 + 12k^2 m^2 - km \\
                &= (7k^3 m^2 + 12k^2 m - k)\cdot m
            \end{align*}
            Since $7k^3 m^2 + 12k^2 m - k$ is an integer, it follows that $m \mid (7n^3+12n^2-n)$.
            \end{proof}
        \item \begin{proof}
            Suppose that $m \mid n$ and $m \mid t$. Then $n = k\cdot m$ and $t = \ell \cdot m$ for
            some integers $k$ and $\ell$. Thus, we have
            \begin{align*}
                n+t = k\cdot m + \ell \cdot m = (k+\ell)\cdot m
            \end{align*}
            Therefore, $m \mid (n+t)$.
            \end{proof}
        \item \begin{proof}
            Suppose $3 \mid 2n$. Then, $2n = k \cdot 3$ for some integer $k$. Combining with the fact
            that $2n$ is even, this implies that $k\cdot 3$ is even. But, $3$ is an odd number, so
            it must be the case that $k$ is even so that $k \cdot 3$ is even. Hence, $k = 2\ell$ for
            some $\ell$. Thus,
            \begin{align*}
                2n = (2\ell) \cdot 3,
            \end{align*}
            and, dividing both sides by $2$, we get
            \begin{align*}
                n = \ell \cdot 3.
            \end{align*}
            Therefore, $3\mid n$.
            \end{proof}
        \item \begin{proof}
            Assume that $m^3 \mid n$ and $n^4 \mid t$. Then, by item (a), $m^6 \mid n^2$, and applying
            it once more, we get $m^{12} \mid n^4$. So, since $n^4 \mid t$, it follows by transitivity
            that $m^{12} \mid t$, as desired.
            \end{proof}
    \end{enumerate}
\end{homeworkProblem}

\begin{homeworkProblem}
    \textit{(Exercise 2.8.)}Prove the following. For each, $m, n,$ and $t$ are integers.
    \begin{multicols}{2}
        \begin{enumerate}[label=(\alph*)]
            \item $1\mid n$.
            \item $n \mid n$.
            \columnbreak
            \item If $mn\mid t$, then $m\mid t$.
            \item If $mn \mid tn$, then $m \mid t$.
        \end{enumerate}
    \end{multicols}
    \textbf{Solution.}
    \begin{enumerate}[label=(\alph*)]
        \item \begin{proof}
            From the identity property for multiplication; that is, $1\cdot n = n\cdot 1 = n$, it
            follows that $1\mid n$.
            \end{proof}
        \item \begin{proof}
            Same as in the proof of item (a). Since $n=1\cdot n$, it follows that $n\mid n$.
            \end{proof}
        \item \begin{proof}
            Suppose $mn \mid t$. Then $t = k\cdot (mn)$ for some $k$. We can rearrange the right hand
            side of the equation to get $t = (kn)\cdot m$, making it more clear that $m \mid t$, as
            desired.
            \end{proof}
        \item \begin{proof}
            Assume that $mn\mid tn$ and that $n\neq 0$. Then, $tn = k\cdot mn$ for some integer $k$.
            Hence, by the cancellation property, it follows that
            \begin{align*}
                t = k\cdot m.
            \end{align*}
            Therefore, $m\mid t$.
            \end{proof}
    \end{enumerate}
\end{homeworkProblem}

\begin{homeworkProblem}
    \textit{(Exercise 2.9.)} Prove that if $m$ and $n$ are positive real numbers and $m < n$, then
    $m^2 < n^2$. You may use the fact that if $a<b$ and $c$ is positive, then $ac<bc$.
    \begin{proof}
        Assume $m$ and $n$ are positive real numbers and that $m < n$. Then, since $m$ is positive, we
        can multiply both sides of the inequality by $m$ and get
        \begin{align*}
            m^2 < mn.
        \end{align*}
        Similarly, since $n$ is positive, we have
        \begin{align*}
            mn < n^2.
        \end{align*}
        Hence, by transitivity,
        \begin{align*}
            m^2 < n^2,
        \end{align*}
        as desired.
    \end{proof}
\end{homeworkProblem}

\begin{homeworkProblem}
    \textit{(Exercise 2.10.)} For each pair of integers, find the unique quotient and remainder when $b$ is divided by $a$.
    \begin{multicols}{3}
        \begin{enumerate}[label=(\alph*)]
            \item $a=4,b=15$
            \item $a=15,b=4$
            \columnbreak
            \item $a=-7,b=16$
            \item $a=5,b=-11$
            \columnbreak
            \item $a=-1,b=-15$
            \item $a=4,b=0$
        \end{enumerate}
    \end{multicols}
    \textbf{Solution.}
    \begin{enumerate}[label=(\alph*)]
        \item $q=3,$ and $r=3$ so that $15 = 4\cdot 3 + 3$
        \item $q=0,$ and $r=4$ so that $4 = 15\cdot 0 + 4$
        \item $q=-2,$ and $r=2$ so that $16 = -7\cdot (-2) + 2$
        \item $q=-2,$ and $r=-1$ so that $-11 = 5\cdot (-2) + (-1)$
        \item $q=15,$ and $r=0$ so that $-15 = (-1)\cdot 15\cdot + 0$
        \item $q=0,$ and $r=0$ so that $0 = 4\cdot 0 + 0$ \qedright
    \end{enumerate}
\end{homeworkProblem}

\begin{homeworkProblem}
    \textit{(Exercise 2.13.)} Define the absolute value of a real number $x$ in this way:
    \begin{align*}
        \left| x \right| = \begin{cases}
            x & \text{if } x\geq 0 \\
            -x & \text{if } x<0.
        \end{cases}
    \end{align*}

    Give three examples showing that if $x$ and $y$ are real numbers, then $|xy|=|x|\cdot|y|$.
    Then, prove that this is true.
    \vspace{5mm}

    \textbf{Solution.}
    If $x=-1,y=3$, then $$|xy|=|-3|=-(-3)=3,$$ and $$|x|\cdot|y|=|-1|\cdot|3|=-(-1)\cdot(3)=3.$$
    If $x=2,y=-3$, then $$|xy|=|-6|=-(-6)=6,$$ and $$|x|\cdot|y|=|2|\cdot|-3|=2\cdot(-(-3))=6.$$
    If $x=-5,y=-4$, then $$|xy|=|20|=20,$$ and $$|x|\cdot|y|=|-5|\cdot|-4|=-(-5)\cdot(-(-4))=5\cdot4=20.$$
    \begin{proof}
        Assume $x$ and $y$ are real numbers. Then, either $x$ and $y$ are both non-negative,
        $x$ and $y$ are both negative, or one of them is negative. We consider these three cases,
        separately. 
        \vspace{5mm}

        \underline{Case 1: $x$ and $y$ are both non-negative.}
        \vspace{5mm}
        Suppose $x\geq 0$ and $y\geq 0$. Then, $|x|=x,|y|=y,$ and $xy\geq 0$. Thus,
        $$|xy|=xy=|x|\cdot|y|.$$

        \underline{Case 2: $x$ and $y$ are both negative.}
        \vspace{5mm}
        Suppose $x<0$ and $y<0$. Then, $|x|=-x, |y|=-y,$ and $xy>0$. Thus, $|xy|=xy$, and
        $$|x|\cdot|y|=(-x)\cdot(-y)=xy.$$
        Therefore, $|xy|=|x|\cdot|y|,$ as desired.
        \vspace{5mm}

        \underline{Case 3: One of $x$ and $y$ is negative.}
        \vspace{5mm}
        Without loss of generality, suppose $x<0$ and $y\geq 0$. Then, $|x|=-x, |y|=y,$ and
        $xy\leq 0$. Hence, $|xy|=-xy,$ and
        $$|x|\cdot|y|=(-x)\cdot y=-xy$$
        Thus, $|xy|=|x|\cdot|y|.$
        \vspace{5mm}

        We have shown that $|xy|=|x|\cdot|y|$ whether $x$ or $y$ is non-negative or negative.
        Combined, this shows that $|xy|=|x|\cdot|y|$ for all real numbers $x,y$.
    \end{proof}
\end{homeworkProblem}

\begin{homeworkProblem}
    \textit{(Exercise 2.14.)} Prove that if $m,n,$ and $t$ are integers, then at least one of $m-n$, $n-t$, and $m-t$
    is even. Also, before your proof, write down three examples of integers $m,n$ and $t$,
    and show which of $m-n,$ $n-t$ or $m-t$ are even.
    \vspace{5mm}

    \textbf{Solution.} If $m=1,n=2,t=3,$ then $m-t$ is even since $m-t=1-3=-2$. If $m=10,n=-1,
    t=7$, then $n-t$ is even since $n-t=(-1)-7=-8$. If $m=4,n=11,t=-8$, then $m-t$ is even
    since $m-t=4-(-8)=12$.
    \vspace{5mm}

    \begin{proof}
        Assume $m,n,$ and $t$ are integers. Then, either $m-n$ is even or $m-n$ is odd. If
        $m-n$ is even, then the result holds. Now, suppose $m-n$ is odd. Then, since
        $$m-n=(m-t)-(n-t),$$
        it follows that $(m-t)-(n-t)$ is odd. But, any two integers must have opposite parities
        if their difference is odd. This implies that either $m-t$ is even or $n-t$ is even.
        So, at least one of $m-n, m-t, n-t$ is even whether $m-n$ is even or $m-n$ is odd.
        Combined, this shows that for any integers $m,n,t$, at least one of $m-n,m-t,n-t$ is
        even.
    \end{proof}
\end{homeworkProblem}

\begin{homeworkProblem}
    \textit{(Exercise 2.15)}
    \begin{enumerate}[label=(\alph*)]
        \item Prove that if $n$ is a positive integer, then $4$ divides $1+(-1)^n(2n-1)$.
        \item Prove that every multiple of $4$ is equal to $1+(-1)^n(2n-1)$ for some positive
        integer $n$.
    \end{enumerate}
    \textbf{Solution.}
    \begin{enumerate}[label=(\alph*)]
        \item Assume $n$ is a positive integer. Then, either $n$ is even or $n$ is odd.
        
        \underline{Case 1: $n$ is even.} Suppose $n$ is even. Then, $n=2k$ for some $k\in
        \mathbb{N}$. Thus,
        \begin{align*}
            1+(-1)^{n}(2n-1) &= 1+(-1)^{2k}(2(2k)-1)\\
            &= 1 + 1\cdot(4k-1) \\
            &= 4k
        \end{align*}
        This implies that, $4$ divides $1+(-1)^n(2n-1)$.
        
        \underline{Case 2: $n$ is odd.} Suppose $n$ is odd. Then, $n=2k+1$ for some $k\in
        \mathbb{N}$. Thus,
        \begin{align*}
            1+(-1)^{n}(2n-1) &= 1+(-1)^{2k+1}(2(2k+1)-1)\\
            &= 1+(-1)(4k+1) \\
            &= 1-4k-1 \\
            &= 4\cdot(-k)
        \end{align*}
        Hence, $4$ divides $1+(-1)^n(2n-1)$. Combined with Case 1, it follows that $4$
        divides $1+(-1)^n(2n-1)$ for any positive integer $n$.\qedright

        \item Let $m$ be a multiple of $4$. Then, $m=4k$ for some $k\in\mathbb{Z}$. Now,
        let $n=2k$. Then,
        $$1+(-1)^{n}(2n-1) = 1+(-1)^{2k}(2(2k)-1).$$
        Note that $(-1)^{2k}=1$ for any $k\in\mathbb{Z}$. So, 
        \begin{align*}
            1+(-1)^{n}(2n-1) &= 1+(2(2k)-1) \\
            &= 1+(4k-1) \\
            &= 4k \\
            &= m,
        \end{align*}
        as desired. \qedright
    \end{enumerate}
\end{homeworkProblem}

\begin{homeworkProblem}
    \textit{(Exercise 2.17)} For each of the following pairs of numbers, list all of their common
    divisors (positive and negative!), and find $\text{gcd}(a,b)$.
    \begin{multicols}{3}
        \begin{enumerate}[label=(\alph*)]
            \item $a=12,b=330$
            \columnbreak
            \item $a=-36,b=64$
            \columnbreak
            \item $a=7,b=-27$
        \end{enumerate}
    \end{multicols}
    \textbf{Solution.}
    \begin{enumerate}[label=(\alph*)]
        \item \underline{Common divisors:} $-6,-3,-2,-1,1,2,3,6$ \\
        $\text{gcd}(12,330)=6$
        \item \underline{Common divisors:} $-4,-2,-1,1,2,4$ \\
        $\text{gcd}(-36,64)=4$
        \item \underline{Common divisors:} $-1,1$ \\
        $\text{gcd}(7,-27)=1$
    \end{enumerate}
\end{homeworkProblem}

\begin{homeworkProblem}
    \textit{(Exercise 2.18.)} For each pair of integers, find gcd$(a,b)$ and integers $k$ and
    $\ell$ such that gcd$(a,b)=ak+b\ell$. (Note: We know that such integers exist by Theorem 2.13.)
    \begin{multicols}{3}
        \begin{enumerate}[label=(\alph*)]
            \item $a=3,b=13$
            \item $a=13,b=3$
            \columnbreak
            \item $a=-25,b=40$
            \item $a=-22,b=-14$
            \columnbreak
            \item $a=62,b=48$
            \item $a=13,b=-50$
        \end{enumerate}
    \end{multicols}
    \begin{enumerate}[label=(\alph*)]
        \item gcd$(3,13)=1$ \\
        Let $k=9$ and $l=-2$. Then,
        $$ak+b\ell=3\cdot9+13\cdot(-2)=27-26=1.$$
        \item gcd$(13,3)=1$ \\
        Let $k=-2$ and $\ell=9$. Then,
        $$ak+b\ell=13\cdot(-2)+3\cdot9=-26+27=1.$$
        \item gcd$(-25,40)=5$ \\
        Let $k=3$ and $\ell=2$. Then,
        $$ak+b\ell=(-25)\cdot3+40\cdot2=-75+80=5.$$
        \item gcd$(-22,-14)=2$ \\
        Let $k=5$ and $\ell=-8$. Then,
        $$ak+b\ell=(-22)\cdot5+(-14)\cdot(-8)=-110+112=2.$$
        \item gcd$(62,48)=2$ \\
        Let $k=7$ and $\ell=-9$. Then,
        $$ak+b\ell=62\cdot7+48\cdot(-9)=434-432=2.$$
        \item gcd$(13,-50)=1$ \\
        Let $k=27$ and $\ell=7$. Then,
        $$ak+b\ell=13\cdot27+(-50)\cdot7=351-350=1.$$
    \end{enumerate}
\end{homeworkProblem}

\begin{homeworkProblem}
    \textit{(Exercise 2.19.)} Let $a$ and $b$ be positive integers, and suppose $r$ is the nonzero
    remainder when $b$ is divided by $a$. Prove that when $-b$ is divided by $a$, the remainder
    is $a-r$.
    \begin{proof}
        Assume that $r\neq 0$. Since $r$ is the remainder when $b$ is divided by $a$, it follows by
        the division algorithm that
        $$b=a\cdot q+r,$$
        where $q\in\mathbb{Z}$ and $0\leq r<a$. Hence,
        $$-b=a\cdot(-q)+(-r).$$
        Now, let $p=-1-q$. Then, $-q=p+1$. So,
        \begin{align*}
            -b&=a\cdot(p+1)+(-r) \\
            &=a\cdot p+(a-r).
        \end{align*}
        Furthermore, $0\leq a-r<a$. So, we have expressed $-b$ in the form that matches with the
        division algorithm. Therefore, by uniqueness, the remainder when $-b$ is divided by $a$
        must be $a-r$.
    \end{proof}
\end{homeworkProblem}

\begin{homeworkProblem}
    \textit{(Exercise 2.20)} Determine the remainder when $3^{302}$ is divided by $28$, and show
    how you found your answer.
    \vspace{1mm}

    \textbf{Solution.} Observe that $3^2\equiv 9 \,(\text{mod }28)$ and $3^6\equiv 1\,(\text{mod }28)$;
    the first one is clear, but the second is true since $728=28\cdot26$, and consequently $728\equiv
    0 \, (\text{mod }28)$. Hence, $(3^6)^{50}\equiv1 \, (\text{mod }28),$ which implies that
    $3^{300}\equiv1 \, (\text{mod }28)$. Thus,
    $$3^{300}\cdot 3^2 \equiv 1\cdot 9 \, (\text{mod }28),$$
    and therefore
    $$3^{302}\equiv 9 \, (\text{mod }28).$$
    So, the remainder when $3^{302}$ is divided by $28$ is $9$.\qedright
\end{homeworkProblem}

\begin{homeworkProblem}
    \textit{(Exercise 2.22.)} Assume that $a$ is an integer and $p$ and $q$ are distinct primes.
    Prove that if $p\mid a$ and $q\mid a$, then $pq\mid a$. Also, before your proof, give three
    examples of this property.
    \vspace{5mm}

    \textbf{Solution.} If $a=12,p=2$ and $q=3$, then $2\mid 12, 3\mid 12$, and $6\mid 12$. If
    $a=70,p=5$ and $q=7$, then $5\mid 70,7\mid 70,$ and $35\mid 70$. If $a=143,p=11$ and $q=13$,
    then $11\mid 143,13\mid 143,$ and $143\mid 143$.

    \begin{proof}
        Suppose $p\mid a$ and $q\mid a$. Then $a=kp$ for some $k\in\mathbb{Z}$, and so
        $q\mid kp$. But, gcd$(p,q)=1$ since $q$ and $p$ are distinct primes. So, by part (ii) of
        Lemma 2.17, $q\mid k$. This implies that $k=\ell q$ for some $\ell\in\mathbb{Z}$. Thus,
        $$a=kp=(\ell q)\cdot p=\ell\cdot(pq),$$
        and therefore $pq\mid a$, as desired.
    \end{proof}
\end{homeworkProblem}

\begin{homeworkProblem}
    \textit{(Exercise 2.23.)} Prove that if $abc$ is a multiple of 10, then at least one of
    $ab,ac,$ or $bc$ is a multiple of 10. Also, before your proof, give three examples of this
    property.
    \vspace{5mm}

    \textbf{Solution.}

    \underline{Examples:}
    \begin{multicols}{3}
        \begin{enumerate}[label=(\alph*)]
            \item $1\cdot 5\cdot 2=10$
            \columnbreak
            \item $5\cdot 5\cdot 8=160$
            \columnbreak
            \item $5\cdot 6\cdot 7=210$
        \end{enumerate}
    \end{multicols}
    \begin{proof}
        Suppose $abc$ is a multiple of $10$. Then $abc$ is also a multiple of $2$ and a
        multiple of $5$. In other words, $5\mid abc$ and $2\mid abc$. Thus, by part (iii) of
        Lemma 2.17., $5\mid ab$ or $5\mid c$, and $2\mid ab$ or $2\mid c$. Equivalently,
        there are three cases:
        \begin{enumerate}[label=(\roman*)]
            \item $5\mid ab$ and $2\mid ab$
            \item One of $ab$ and $c$ are divisible by $5$, and the other is divisible by $2$
            \item $5\mid c$ and $2\mid c$
        \end{enumerate}
        It is clear to see that the result holds for cases (i) and (iii) --- For case (i),
        it follows by the previous exercise that $10\mid ab$; while, for case (iii), we have
        $10\mid c$ which implies that $10\mid ac$ and $10\mid bc$. Now, we'll show that the
        result also holds for case (ii).
        \vspace{3mm}

        \underline{Case (ii): $5\mid ab$ and $2\mid c$.} Without loss of generality,
        suppose $5\mid ab$ and $2\mid c$. Then, by part (iii) of Lemma 2.17., either
        $5\mid a$ or $5\mid b$. So, considering that $2\mid c$, it must be the case that
        either $10\mid ac$ or $10\mid ac$.
        \vspace{3mm}

        Combining all these three cases, it follows that at least one of $ab,ac,$ or $bc$
        is a multiple of $10$.
    \end{proof}
\end{homeworkProblem}
\end{document}